% =============================================================================
% APPENDIX CK: FORMAL-BELIEFOPPOSITIONS: Social Democratic vs. Liberal vs. Conservative Belief Contrasts
% =============================================================================
% Module: BCM2_FORMAL_CONTRAST (Ideological Comparison)
% Category: FORMAL
% Version: 1.0 (January 2026)
% Status: Final
% Language: English
%
% This appendix maps the 15 social democratic beliefs (from Appendix CJ) against
% liberal and conservative alternatives, showing where logical compromise is possible
% (Secondary Beliefs) and where it is not (Deep Core / Policy Core).
%
% =============================================================================

\begin{tcolorbox}[colback=purple!5!white, colframe=purple!75!black,
    title=Chapter Linkage]

\textbf{Primary Chapter:} Chapter 14: Application - Democratic Resilience

\textbf{Related Appendices:}
\begin{itemize}[nosep]
\item Appendix CJ (FORMAL-BELIEFARCHITECTURE): Sabatier ACF framework
\item Appendix BF (DOMAIN-SOCIALDEMOCRACY): 10 normative principles
\end{itemize}

\end{tcolorbox}

% =============================================================================

\chapter{FORMAL-BELIEFOPPOSITIONS: Ideological Contrasts in Deep, Policy, and Secondary Beliefs}
\label{app:ck-beliefoppositions}


\begin{tcolorbox}[colback=blue!5!white, colframe=blue!75!black,
    title=Quick Reference: Appendix CK]

\textbf{Appendix:} CK (FORMAL)

\textbf{Core Question:} What are the key findings in this domain?

\textbf{Key Concepts:}
\begin{itemize}[nosep]
\item Framework integration
\item Parameter validation
\item Cross-references
\end{itemize}

\textbf{Cross-References:} See related appendices below
\end{tcolorbox}

\begin{abstract}

This appendix applies comparative ideological analysis to the 15 social democratic beliefs, showing how liberalism and conservatism propose different answers to the same fundamental questions. The analysis demonstrates that while Secondary Beliefs (concrete policies) allow compromise, Policy Core and especially Deep Core beliefs create logical incompatibilities that make coalition-building difficult. The Austrian case shows how grand coalitions (SPÖ/ÖVP) can survive when Deep Core differences are muted (``We both value democracy''), but break when Policy Core is violated (Greens demand media regulation; ÖVP refuses).

\end{abstract}

\begin{tcolorbox}[colback=yellow!5!white, colframe=yellow!75!black,
    title=Quick Reference: The Three-Ideology Matrix}

\begin{table}[h]
\centering
\footnotesize
\renewcommand{\arraystretch}{1.2}
\begin{tabular}{@{}lccc@{}}
\toprule
\textbf{Question} & \textbf{Social Democratic} & \textbf{Liberal} & \textbf{Conservative} \\
\midrule
\textbf{Deep Core:} & & & \\
Equal worth? & Yes, intrinsic & Yes, procedural & Hierarchical roles \\
Inequality ok? & No (problematic) & If earned & Natural/incentive \\
Solidarity? & Obligatory & Voluntary & Family/tradition \\
Material→Freedom? & Both required & Formal rights enough & Virtue/duty \\
Society malleable? & Yes, reform possible & Yes, but bounded & Caution/organic \\
\midrule
\textbf{Policy Core:} & & & \\
State responsible? & Yes, comprehensive & Minimal, subsidiarity & Traditional inst. \\
Opportunity? & Active intervention & Equal starting line & Merit/inheritance \\
Markets regulated? & Yes, systematically & Minimal/competition & Traditional limits \\
Redistribution legit? & Yes, systematic & No, violates property & Charity/noblesse \\
Labor special? & Yes, powerful unions & Market price OK & Employer authority \\
\midrule
\textbf{Secondary:} & & & \\
Cash transfers? & Universal, generous & Means-tested & Charity/family \\
Public services? & Universal, public & Subsidized choice & Private/traditional \\
Regulation? & Extensive & Minimal exceptions & Customary limits \\
Progressive tax? & High rates (50\%+) & Flat/minimal & Lower rates \\
Codetermination? & Yes, mandatory & No, contract & Employer-led \\
\bottomrule
\end{tabular}
\end{table}

\textbf{Compatibility Assessment:} SD ↔ Liberal = High (negotiate on Secondary). SD ↔ Conservative = Low (clash on Policy Core). Liberal ↔ Conservative = Medium (both accept private property).

\end{tcolorbox}

% =============================================================================
% SECTION 1: FUNDAMENTAL QUESTION & AXIOMS
% =============================================================================

% -----------------------------------------------------------------------------
% APPENDIX SCOPE BOX
% -----------------------------------------------------------------------------

\begin{tcolorbox}[
    colback=orange!5!white,
    colframe=orange!75!black,
    title={\textbf{Appendix Scope: Ziel / In-Scope / Out-of-Scope / Constraints}},
    fonttitle=\bfseries
]

\textbf{Ziel (Objective):}
\begin{itemize}[nosep]
    \item How do Deep Core ideological differences constrain the feasible set of coalition policies, and when do they become salient enough to destabilize coalitions?
\end{itemize}

\textbf{In-Scope:}
\begin{itemize}[nosep]
    \item The Fundamental Question
    \item Axioms and Formal Definitions
    \item Deep Core Beliefs: Fundamental Incompatibilities
    \item Policy Core Beliefs: Limited Negotiation Space
    \item Secondary Beliefs: Maximum Negotiation Space
\end{itemize}

\textbf{Out-of-Scope (delegiert an andere Appendices):}
\begin{itemize}[nosep]
    \item FORMAL-BELIEFARCHITECTURE $\rightarrow$ Appendix CJ
    \item DOMAIN-SOCIALDEMOCRACY $\rightarrow$ Appendix BF
    \item FORMAL-BELIEFARCHITECTURE $\rightarrow$ Appendix CJ
    \item DOMAIN-SOCIALDEMOCRACY $\rightarrow$ Appendix BF
    \item CONTEXT-SPÖ_PARADOX $\rightarrow$ Appendix BH
\end{itemize}

\textbf{Constraints (Anwendungsgrenzen):}
\begin{itemize}[nosep]
    \item [TODO: Define when this appendix does NOT apply]
    \item [TODO: Define required prerequisites]
    \item [TODO: Define known limitations]
\end{itemize}

\textbf{Lieferobjekte (Deliverables):}
\begin{itemize}[nosep]
    \item 8 Table(s)
    \item 1 Equation(s)
    \item 2 Axiom(s)
    \item Worked Example(s)
\end{itemize}

\end{tcolorbox}


\section{The Fundamental Question}
\label{sec:ck-fundamental}

\begin{tcolorbox}[colback=red!5!white, colframe=red!75!black,
    title=The Fundamental Question]

\textbf{How do Deep Core ideological differences constrain the feasible set of coalition policies, and when do they become salient enough to destabilize coalitions?}

\end{tcolorbox}

% =============================================================================

\section{Axioms and Formal Definitions}
\label{sec:ck-axioms}

\begin{tcolorbox}[colback=gray!5!white, colframe=gray!75!black,
    title=Axiom CK-1: Coalition Possibility Frontier}

\textbf{Statement:} $\text{CPF} = \{ \text{Policies} : \text{Satisfy all partners' Policy Core}\}$

\textbf{Interpretation:} Feasible coalition policies must lie within the intersection of all coalition partners' core constraints.

\end{tcolorbox}

\begin{tcolorbox}[colback=gray!5!white, colframe=gray!75!black,
    title=Axiom CK-2: Deep Core as Hard Constraint}

\textbf{Statement:} Deep Core beliefs are non-negotiable; they define coalition membership.

\textbf{Implication:} When Deep Core differences become salient (not suppressed), coalitions collapse.

\end{tcolorbox}

% =============================================================================

\section{Deep Core Beliefs: Fundamental Incompatibilities}
\label{sec:ck-deepcore}

\subsection{Deep Core 1: Equal Moral Worth}

\begin{tcolorbox}[colback=red!10!white, colframe=red!75!black,
    title=Belief 1: Human Equality]

\textbf{Social Democratic Position:}
\begin{quote}
``All humans possess intrinsic, equal worth. This is not earned; it is fundamental. Therefore, extreme inequality violates basic dignity.''
\end{quote}

\textbf{Liberal Position:}
\begin{quote}
``All humans have equal formal/procedural worth (equal before law). But in markets, unequal outcomes reflect unequal effort, talent, and contribution. This is just if procedures are fair.''
\end{quote}

\textbf{Conservative Position:}
\begin{quote}
``Humans have different capacities, roles, and virtues. Society requires hierarchy and complementary stations. Equality is impossible and undesirable; stability through differentiation is better.''
\end{quote}

\textbf{Logical Compatibility:}

\begin{table}[h]
\centering
\small
\renewcommand{\arraystretch}{1.2}
\begin{tabular}{@{}lll@{}}
\toprule
& \textbf{SD vs. Liberal} & \textbf{SD vs. Conservative} \\
\midrule
Compatibility & Medium & Low \\
Negotiation Space & Outcomes (SD flexible on income inequality if opportunity equal) & Ideology (fundamentally different) \\
Coalition Basis & Economic justice + procedural fairness & Unlikely without suppression \\
\bottomrule
\end{tabular}
\end{table}

\textbf{Austrian Example:} SPÖ/ÖVP grand coalitions work because both parties suppress Deep Core differences (emphasize ``shared democracy'') and focus on Secondary Belief negotiation. But when policy coherence breaks (ÖVP protects media concentration; SPÖ rhetorically opposes it), Deep Core differences become salient.

\end{tcolorbox}

\subsection{Deep Core 2: Inequality as Problem}

\begin{tcolorbox}[colback=red!10!white, colframe=red!75!black,
    title=Belief 2: Inequality Problem}

\textbf{Social Democratic:} ``Inequality beyond a threshold is morally indefensible.''

\textbf{Liberal:} ``Inequality resulting from fair procedures is acceptable. Concern is poverty, not inequality per se.''

\textbf{Conservative:} ``Inequality is natural, incentive-generating, and compatible with mutual obligation (noblesse oblige).''

\textbf{Logical Conflict:} These are \textbf{not empirically resolvable}. They are axiomatic positions about justice. No amount of data on inequality's effects can shift a Deep Core belief.

\textbf{Political Consequence:} SD/Liberal coalitions can compromise on ``how much redistribution'' (negotiable). SD/Conservative coalitions must either ignore the disagreement (as Austria did) or fight on principle.

\end{tcolorbox}

\subsection{Deep Core 3: Solidarity as Obligation}

\begin{tcolorbox}[colback=red!10!white, colframe=red!75!black,
    title=Belief 3: Solidarity}

\textbf{Social Democratic:} ``Members of a political community \textit{owe} each other support. Mutual obligation is constitutive.''

\textbf{Liberal:} ``Individuals are atomistic; solidarity is voluntary (charity). Obligations are contractual, not inherited.''

\textbf{Conservative:} ``Solidarity flows through traditional institutions (family, church, local community), not the state. State-imposed solidarity is illegitimate.''

\textbf{Incompatibility:} These differ on the \textbf{source} of obligation (intrinsic vs. contractual vs. traditional). Compromise is superficial.

\textbf{Austria's Problem (2000-2015):} SPÖ maintained ``We are solidaristic'' while ÖVP maintained ``Solidarity is optional.'' On pensions (accepted social insurance history), both could work together. On media concentration (new problem), neither activated their supposed principle.

\end{tcolorbox}

% =============================================================================

\section{Policy Core Beliefs: Limited Negotiation Space}
\label{sec:ck-policycore}

\subsection{Policy Core 6: State Responsibility for Social Security}

\begin{tcolorbox}[colback=blue!10!white, colframe=blue!75!black,
    title=Belief 6: Welfare State}

\textbf{Social Democratic:}
\begin{quote}
``State is primarily responsible for universal social insurance (health, unemployment, pensions). Coverage must be comprehensive and non-stigmatizing.''
\end{quote}

\textbf{Liberal:}
\begin{quote}
``State provides safety net for the destitute (means-tested). Most people should insure privately or through employers. Universalism creates moral hazard.''
\end{quote}

\textbf{Conservative:}
\begin{quote}
``Families and charities are primary. State involvement should be minimal and temporary (emergency relief). Permanent state programs undermine character.''
\end{quote}

\textbf{Negotiation Possibilities (Policy Core level):}

\begin{table}[h]
\centering
\small
\renewcommand{\arraystretch}{1.3}
\begin{tabular}{@{}p{3cm}ccc@{}}
\toprule
\textbf{Issue} & \textbf{SD Position} & \textbf{Liberal Deal} & \textbf{Conservative Deal} \\
\midrule
Coverage & Universal (all) & Targeted (poor) & Selective (deserving) \\
Funding & Progressive tax & User fees + tax & Voluntary charity \\
Generosity & High (80\% replacement) & Moderate (50\%) & Minimal (poverty line) \\
Public vs. Private & Public primary & Mixed & Private primary \\
\bottomrule
\end{tabular}
\end{table}

\textbf{Austrian Coalition Practice:} SPÖ/ÖVP compromised at SD/Liberal interface: Universal coverage (SD wins) + some private options (Liberal wins). Not full SD universalism, not full Conservative minimalism.

\textbf{Why Compromise Broke in Media (2000-2015):} Media concentration was not yet an established social insurance issue. SPÖ argued ``Market will fix it,'' ÖVP agreed. Result: No regulation at all (neither SD nor Liberal nor Conservative position consistently applied).

\end{tcolorbox}

\subsection{Policy Core 8: Markets Require Regulation}

\begin{tcolorbox}[colback=blue!10!white, colframe=blue!75!black,
    title=Belief 8: Market Regulation}

\textbf{Social Democratic:}
\begin{quote}
``Markets systematically concentrate power and produce injustice. Comprehensive regulation (labor, finance, media, land, etc.) is necessary and legitimate.''
\end{quote}

\textbf{Liberal:}
\begin{quote}
``Competition prevents power concentration. Regulation should target specific market failures (monopoly, information asymmetry). But default should be competition.''
\end{quote}

\textbf{Conservative:}
\begin{quote}
``Markets embody traditional property rights and custom. Regulation should preserve these, not redistribute. Heavy-handed intervention is destructive.''
\end{quote}

\textbf{Critical Negotiation Point:}

\begin{table}[h]
\centering
\small
\renewcommand{\arraystretch}{1.3}
\begin{tabular}{@{}p{3.5cm}ccc@{}}
\toprule
\textbf{Market Sector} & \textbf{SD Default} & \textbf{Liberal Agreement} & \textbf{Conservative Resistance} \\
\midrule
Labor markets & ✅ Regulate (strong) & ✅ OK (market failure) & ❌ No (union power bad) \\
Finance markets & ✅ Regulate (systemic risk) & ✅ OK (too-big-to-fail) & ❌ No (limits profit) \\
Consumer protection & ✅ Regulate (asymmetry) & ✅ OK (information) & ⚠️ Maybe (light touch) \\
Media markets & ✅ Regulate (democracy) & ⚠️ Hesitant (free speech) & ❌ No (state control risk) \\
Land markets & ✅ Regulate (inequality) & ⚠️ Hesitant (property) & ❌ No (violates property) \\
\bottomrule
\end{tabular}
\end{table}

\textbf{Austria's Fatal Incoherence:} SPÖ regulated labor/finance (Liberal-SD compromise). But SPÖ \textit{rhetoric} said ``We regulate markets'' (Policy Core 8) while \textit{practice} ignored media/land (yielded to Conservative resistance). This inconsistency -- regulation sometimes, sometimes not -- created coherence perception of hypocrisy.

\textbf{The Sabatier Prediction:} Voters detected: ``You claim to regulate markets for democracy, but Krone dominates unchecked. You're lying about your principles.'' Defection to FPÖ (which at least said ``Elites corrupt'' without pretending to fix it).

\end{tcolorbox}

\subsection{Policy Core 9: Redistribution Is Legitimate}

\begin{tcolorbox}[colback=blue!10!white, colframe=blue!75!black,
    title=Belief 9: Redistribution}

\textbf{Social Democratic:}
\begin{quote}
``Systematic redistribution (progressive tax + transfers) is ethically necessary and economically productive.''
\end{quote}

\textbf{Liberal:}
\begin{quote}
``Limited redistribution to prevent destitution is acceptable. But beyond that, high redistribution creates moral hazard and efficiency loss.''
\end{quote}

\textbf{Conservative:}
\begin{quote}
``Redistribution violates property rights and the incentive system. Charity may be appropriate, but not coercive taxation.''
\end{quote}

\textbf{Possible Compromise (SD/Liberal):}

Progressive taxation up to a certain level (e.g., top marginal rate 45-50\%) is negotiable. Beyond that, SD and Liberal diverge significantly. Conservative refuses progressive taxation entirely.

\textbf{Austria's Implementation:}
\begin{itemize}
\item Top income tax: 55\% (SD wins)
\item Wealth tax: Abolished 2000 (Conservative wins)
\item Result: Partially coherent with Policy Core 9, but weakened by wealth tax loss
\end{itemize}

\end{tcolorbox}

% =============================================================================

\section{Secondary Beliefs: Maximum Negotiation Space}
\label{sec:ck-secondary}

\subsection{Secondary Belief 11: Cash Transfers}

\begin{tcolorbox}[colback=green!10!white, colframe=green!75!black,
    title=Belief 11: Cash Transfer Design}

\textbf{Social Democratic:} Universal child benefits (no means-test).

\textbf{Liberal:} Targeted to low-income families (means-test, but automatic).

\textbf{Conservative:} Voluntary charity; family responsibility.

\textbf{Room for Compromise:}

All three can implement ``child benefits,'' but disagree on:
- Universality (SD: yes, Liberal: no, Conservative: no)
- Generosity level (SD: high, Liberal: medium, Conservative: minimal)
- Delivery (SD: public, Liberal: public/private mix, Conservative: private charity)

\textbf{Austrian Practice:} Adopted hybrid: Universal eligibility + modest level (~€160/month) + public delivery. Satisfies SD principle (universal) + Liberal concern (not too generous) + Conservative tradition (family responsibility acknowledged).

\textbf{Testability:} If research shows means-testing reduces stigma, all three would consider adjustment. This is Secondary Belief flexibility.

\end{tcolorbox}

\subsection{Secondary Belief 13: Regulatory Intervention}

\begin{tcolorbox}[colback=green!10!white, colframe=green!75!black,
    title=Belief 13: What to Regulate (Maximum Conflict)}

\textbf{Labor Markets:}
\begin{itemize}
\item SD: Comprehensive (min wage, hours, safety, collective bargaining)
\item Liberal: Minimal (only monopoly power, then OK)
\item Conservative: No (union power too much)
\item \textbf{Austrian Outcome:} Strong labor regulation (SD/Liberal compromise held)
\end{itemize}

\textbf{Finance Markets:}
\begin{itemize}
\item SD: Comprehensive (capital controls, speculation taxes)
\item Liberal: Focused (too-big-to-fail, transparency)
\item Conservative: No (limits profit)
\item \textbf{Austrian Outcome:} Moderate regulation (post-2008, Liberal/SD aligned)
\end{itemize}

\textbf{Media Markets:}
\begin{itemize}
\item SD: Ownership caps, transparency, public broadcaster support
\item Liberal: Hesitant (free speech concerns); maybe ownership caps
\item Conservative: No (state control risk)
\item \textbf{Austrian Outcome:} NO regulation (Conservative veto held; Liberal/SD split)
\end{itemize}

\textbf{The Divergence:} Labor and finance regulation allowed SD/Liberal compromise because the market failure argument (wage suppression, systemic risk) was empirically persuasive. Media regulation failed because:
- SD saw ``ownership concentration'' as market failure
- Liberal worried about ``state control of news''
- Conservative rejected it entirely
- Without clear empirical consensus, ÖVP (Conservative) vetoed

\textbf{Babler's Genius:} Reframe media regulation as ``transparency + competition'' (Liberal language), not ``state control'' (Conservative fear). Marketing, but logically valid.

\end{tcolorbox}

\subsection{Secondary Belief 14: Progressive Taxation}

\begin{tcolorbox}[colback=green!10!white, colframe=green!75!black,
    title=Belief 14: Tax Rate Levels (Maximum Compromise)}

\begin{table}[h]
\centering
\small
\renewcommand{\arraystretch}{1.3}
\begin{tabular}{@{}lcccc@{}}
\toprule
\textbf{Level} & \textbf{SD Target} & \textbf{Liberal Compromise} & \textbf{Conservative Limit} & \textbf{Austria 2026} \\
\midrule
Top income rate & 60-65\% & 45-50\% & 30\% & 55\% \\
Wealth tax & Yes, 1-2\% & Maybe, <1\% & No & 0\% (abolished) \\
Inheritance tax & Yes, 40\% & Maybe, 20\% & No & ~7-14\% \\
\bottomrule
\end{tabular}
\end{table}

\textbf{Why These Are Secondary (Not Policy Core):}

The \textbf{principle} (progressive taxation) is Policy Core. But the \textbf{exact rates} depend on empirical evidence (revenue, investment behavior, growth effects).

\textbf{Negotiation Typical in Austria:} SPÖ proposed 60\% top rate; ÖVP rejected as ``excessive.'' They compromised at 55\%. Both could claim victory: SPÖ got progression, ÖVP limited it.

\textbf{Where Compromise Failed:} Wealth tax. SPÖ/ÖVP coalition abolished it (2000) under pressure from capital interests. Neither party maintained Policy Core 9 (redistribution legitimate) when capital interests pushed.

\textbf{Lesson:} Secondary Beliefs are flexible, but only if political will exists. If capital interests lobby successfully against wealth tax, and neither SD nor Liberals resist, then the redistributive principle erodes.

\end{tcolorbox}

% =============================================================================

\section{The Compromise Frontier: What Can Coalition Partners Negotiate?}
\label{sec:ck-frontier}

\subsection{Coalition Possibility Frontier (CPF)}

Not all policy combinations are feasible. The set of feasible policies forms a \textbf{Coalition Possibility Frontier} defined by the intersection of parties' ideological constraints.

\begin{equation}
\text{CPF} = \{ \text{Policies} : \text{Satisfy all coalition partners' Policy Core}\}
\end{equation}

\textbf{Example: SPÖ/ÖVP Feasible Set (Austria)}

\begin{table}[h]
\centering
\small
\renewcommand{\arraystretch}{1.3}
\begin{tabular}{@{}lcc@{}}
\toprule
\textbf{Policy Domain} & \textbf{Inside CPF?} & \textbf{Example} \\
\midrule
Pensions & ✅ Yes & Both support PAYG system \\
Public health & ✅ Yes & Both support universal insurance \\
Education & ✅ Yes & Both support public schools \\
Labor regulation & ✅ Yes & Both support works councils \\
Progressive tax & ⚠️ Marginal & ~50\% top rate (ÖVP limit) \\
Wealth tax & ❌ No & ÖVP refuses entirely \\
Media regulation & ❌ No & ÖVP refuses media ownership caps \\
Environmental reg & ⚠️ Marginal & Greens push; ÖVP resists \\
\bottomrule
\end{tabular}
\end{table}

\textbf{Critical Finding:} Media regulation fell \textbf{outside} the SPÖ/ÖVP CPF because:
- SPÖ Policy Core 8: ``Markets need regulation''
- ÖVP Policy Core 8: ``Markets need regulation (except media, which is property-based)''

The parties did not actually disagree on the principle. ÖVP simply prioritized Conservative \textit{Deep Core} (property rights, hierarchy) over the shared Policy Core (market regulation for democracy).

When Babler offered \textit{partial} media regulation (not full nationalization), it could have entered the CPF. But 2016-2020 ÖVP refused anyway.

\subsection{Why 2016 Media Reform Failed (CPF Analysis)}

\begin{table}[h]
\centering
\small
\renewcommand{\arraystretch}{1.4}
\begin{tabular}{@{}lccc@{}}
\toprule
\textbf{Party} & \textbf{Media Reg Position} & \textbf{Policy Core} & \textbf{Deep Core Conflict} \\
\midrule
SPÖ & Yes (ownership caps) & Market regulation necessary & Doesn't prioritize property \\
Greens & Yes (transparency + caps) & Market regulation necessary & Doesn't prioritize property \\
ÖVP & No (market handles) & Market regulation (selective) & Property rights sacred \\
\bottomrule
\end{tabular}
\end{table}

\textbf{CPF Analysis:}
\begin{itemize}
\item SPÖ/Greens coalition could implement media regulation (Policy Core aligned)
\item But ÖVP was also in government (2020-2024)
\item ÖVP's Deep Core (property → hierarchy → media ownership unregulated) contradicted CPF point
\item Result: Media regulation remained outside the three-party CPF
\end{itemize}

\textbf{Babler's Strategic Move (2025):} Govern as SPÖ-led, without ÖVP. This shifts CPF: SPÖ/Greens alone can implement media regulation. By excluding ÖVP's Deep Core veto, the policy enters the feasible set.

\end{tcolorbox}

% =============================================================================

\section{Key Results: Belief Incompatibility Predicts Coalition Instability}
\label{sec:ck-results}

\begin{tcolorbox}[colback=green!5!white, colframe=green!75!black,
    title=Result CK-R1: Deep Core Incompatibility Requires Suppression}

\textbf{Finding:} Grand coalitions (SPÖ/ÖVP) can coexist only by \textbf{not discussing} Deep Core differences (equality, hierarchy, property). Silence on fundamental values allows policy compromises.

\textbf{Austrian Evidence:}
\begin{itemize}
\item Parties agree to focus on ``practical matters'' (pensions, education, labor)
\item Parties do NOT debate ``Is hierarchy natural?'' or ``Is property sacred?''
\item Result: Coalition stability on Secondary/some Policy Core issues
\item Fragility: When new issues arise (media) that implicate suppressed Deep Core, coalition breaks
\end{itemize}

\end{tcolorbox}

\begin{tcolorbox}[colback=green!5!white, colframe=green!75!black,
    title=Result CK-R2: Policy Core Violation → Coalition Collapse}

\textbf{Finding:} When one coalition partner violates its own Policy Core while the other upholds theirs, the coalition destabilizes.

\textbf{Austrian Example (2000-2015):}
\begin{itemize}
\item SPÖ stated Policy Core 8: ``Markets need regulation''
\item ÖVP stated Policy Core 8 (selective): ``Some markets need regulation''
\item Reality: Media never regulated (violated SPÖ's Policy Core 8)
\item SPÖ accepted this, violating its own coherence (Appendix CJ finding)
\item Voters detected incoherence → SPÖ support collapsed
\end{itemize}

\end{tcolorbox}

\begin{tcolorbox}[colback=green!5!white, colframe=green!75!black,
    title=Result CK-R3: Secondary Belief Compromise Is Stable}

\textbf{Finding:} On Secondary Beliefs (concrete policy instruments), all three ideologies can negotiate because the empirical evidence is falsifiable.

\textbf{Examples:}
\begin{itemize}
\item Progressive tax rate: 50\% is compromise between SD (60\%) and Conservative (30\%). All three accept.
\item Child benefits: Universal + modest level is compromise. All three accept.
\item Labor regulation: Workday limits, minimum wages, safety rules all negotiable at margins.
\end{itemize}

\textbf{Stability:} These compromises hold until evidence shifts (e.g., if research shows high top tax rates reduce growth, all three might negotiate downward).

\end{tcolorbox}

% =============================================================================

\section{Summary}
\label{sec:ck-summary}

\begin{tcolorbox}[colback=green!10!white, colframe=green!75!black,
    title=Appendix CK: Summary]

\textbf{Core Contribution:}
\begin{itemize}[nosep]
    \item Maps 15 social democratic beliefs against liberal and conservative alternatives
    \item Shows where compromise is possible (Secondary Beliefs) and where it is not (Deep/Policy Core)
    \item Demonstrates that Austria's grand coalition worked by suppressing Deep Core differences
\end{itemize}

\textbf{Key Insights:}
\begin{itemize}[nosep]
    \item Deep Core incompatibility requires coalition partners to maintain ``constructive silence''
    \item Policy Core violations (like media regulation) signal incoherence that voters detect
    \item Secondary Beliefs allow maximum compromise space
    \item Coalition Possibility Frontier (CPF) shrinks as Deep Core differences become salient
\end{itemize}

\end{tcolorbox}

% =============================================================================

\section{Critical Foundations and Objections}
\label{sec:ck-foundations}

\subsection{Objection 1: Are Deep Core Beliefs Actually Non-Negotiable?}

\textbf{Objection:} ``You claim Deep Core beliefs are immutable, but people change their minds all the time. Core beliefs do shift.''

\textbf{Response:} Sabatier distinguishes between \textit{individual} belief change (yes, happens) and \textit{coalition-level} belief change (rare, generational). A single social democrat might come to believe ``hierarchy is natural.'' But for that to reshape the social democratic coalition's Deep Core would take decades and replacement of multiple generations of leadership.

\subsection{Objection 2: Isn't the Three-Party Framework Too Rigid?}

\textbf{Objection:} ``Real politics has more than three ideologies. There are anarchists, monarchists, theocrats, libertarians. Your SD/Liberal/Conservative triadic oversimplifies.''

\textbf{Response:} This appendix focuses on the Overton window of Austrian politics (1970-2025). Within that window, the three frameworks reasonably cover the feasible ideological space. For other countries or historical periods, the framework may need adjustment.

% =============================================================================

\section{Open Issues}
\label{sec:ck-open}

\begin{enumerate}

\item \textbf{Can Deep Core differences ever be explicitly negotiated?} Austria suppressed them (silence = stability). But is there a stable equilibrium where parties \textit{disagree openly} on Deep Core values and still coalesce? Test: Compare grand coalitions across EU (Germany, France, Belgium).

\item \textbf{How much Policy Core incoherence can a party tolerate?} SPÖ violated Policy Core 8 (media regulation) for 15 years. How long can a party maintain voter support while violating its own Policy Core? Testable threshold?

\item \textbf{Does the CPF shrink over time?} As Deep Core differences become more salient (migration, climate, tech), do coalition partners find fewer policies within the feasible set? Predict: European grand coalitions become less stable 2020-2030.

\item \textbf{Can media regulation be re-framed to satisfy all ideologies?} Babler is trying: ``Transparency regulation'' (Liberal language, SD purpose). If it works, this suggests Deep Core conflicts can be managed through rhetorical re-framing. But is the disagreement only rhetorical, or truly fundamental?

\end{enumerate}

% =============================================================================

\section{References}
\label{sec:ck-references}

\subsection{Ideological Analysis}

\begin{itemize}[nosep]
\item Leftwich, A. (1984). \textit{What is Politics? The Activity and its Study}. Oxford University Press.

\item Heywood, A. (2007). \textit{Political Ideologies: An Introduction} (4th ed.). Palgrave Macmillan.

\item Freeden, M. (2003). \textit{Ideology: A Very Short Introduction}. Oxford University Press.

\end{itemize}

\subsection{Coalition Theory}

\begin{itemize}[nosep]
\item Riker, W. H. (1962). \textit{The Theory of Political Coalitions}. Yale University Press.

\item Laver, M., \& Shepsle, K. A. (1996). \textit{Making and Breaking Governments: Cabinets and Legislatures in Parliamentary Democracies}. Cambridge University Press.

\end{itemize}

\nocite{bcm_master}

\subsection{General References}

For comprehensive references on ideological analysis, coalition theory, and belief systems, consult the master bibliography.

\subsection{Glossary of Symbols and Terms}

\begin{tcolorbox}[colback=blue!5!white, colframe=blue!75!black]
\textbf{Central Reference:} For complete symbol definitions, see
\textbf{Appendix GLS} (\textit{Glossary and Notation}).
\end{tcolorbox}

\begin{table}[h]
\centering
\small
\renewcommand{\arraystretch}{1.3}
\begin{tabular}{@{}llp{6cm}@{}}
\toprule
\textbf{Term} & \textbf{Name} & \textbf{Definition} \\
\midrule
Deep Core & Fundamental beliefs & Non-negotiable axioms (moral, epistemological) \\
Policy Core & Stable commitments & Broad policy orientations; negotiable at margins \\
Secondary & Concrete policies & Specific instruments; highly falsifiable \\
CPF & Coalition Possibility Frontier & Set of feasible policies satisfying all coalition partners \\
Coherence & Internal consistency & Alignment between stated Policy Core and actual Secondary implementation \\
Incoherence & Inconsistency & Policy Core violated by Secondary implementation \\
\bottomrule
\end{tabular}
\end{table}

% =============================================================================

\begin{tcolorbox}[colback=orange!5!white, colframe=orange!75!black,
    title=Cross-Reference Links]

\textbf{Related Appendices:}
\begin{itemize}[nosep]
\item Appendix CJ (FORMAL-BELIEFARCHITECTURE): Sabatier ACF framework applied to SD
\item Appendix BF (DOMAIN-SOCIALDEMOCRACY): The 10 normative principles
\item Appendix BH (CONTEXT-SPÖ_PARADOX): How ÖVP veto paralyzed SPÖ reform
\end{itemize}

\end{tcolorbox}

% =============================================================================
\end{document}
